\documentclass[11pt]{article}
\usepackage[margin=0.75in]{geometry}
\usepackage{fancyhdr}
\usepackage{array}
\usepackage{booktabs}
\usepackage{xcolor}
\usepackage{enumitem}

\pagestyle{fancy}
\fancyhf{}
\rhead{IME Report - Santos, Maria E.}
\lhead{Pacific Medical Evaluators}
\rfoot{Page \thepage}

\definecolor{imeblue}{RGB}{0,51,102}

\begin{document}

\begin{center}
{\Large\bfseries\color{imeblue} PACIFIC MEDICAL EVALUATORS, INC.}\\[0.3em]
{\small Independent Medical Examination Services}\\
{\small 2000 Professional Tower, Suite 1500, Los Angeles, CA 90017}\\
{\small Phone: (213) 555-9000}\\[1em]
{\LARGE\bfseries INDEPENDENT MEDICAL EXAMINATION REPORT}
\end{center}

\vspace{0.5em}
\hrule height 2pt
\vspace{1em}

\section*{EXAMINATION INFORMATION}

\begin{tabular}{@{}p{4cm}p{10cm}@{}}
\textbf{Examinee:} & Maria Elena Santos \\
\textbf{Date of Birth:} & March 14, 1988 \\
\textbf{Date of Examination:} & April 15, 2026 \\
\textbf{Date of Injury:} & February 8, 2026 \\
\textbf{Examiner:} & Richard Stevenson, MD \\
\textbf{Specialty:} & Orthopedic Surgery \\
\textbf{Requesting Party:} & Hartford Insurance Company \\
\textbf{Claim Number:} & HRT-2026-0892147 \\
\end{tabular}

\section*{PURPOSE OF EXAMINATION}

This independent medical examination was requested by Hartford Insurance Company to address the following:

\begin{enumerate}
\item Current diagnosis and medical status
\item Causation of current conditions
\item Maximum Medical Improvement (MMI) status
\item Permanent impairment rating, if applicable
\item Appropriateness of current treatment
\item Work capacity and restrictions
\item Future medical needs
\end{enumerate}

\section*{RECORDS REVIEWED}

\begin{itemize}[nosep]
\item Riverside General Hospital ER records (02/08/2026)
\item Dr. Kevin Park orthopedic records (02/10/2026, 03/08/2026)
\item Dr. Sarah Goldstein neurology consultation (02/12/2026)
\item Riverside Imaging Center MRI report (02/15/2026)
\item ProMotion Physical Therapy records (02/19/2026 - 04/01/2026)
\item Employer incident report
\item OSHA investigation report
\end{itemize}

\section*{HISTORY AS OBTAINED FROM EXAMINEE}

Ms. Santos is a 35-year-old right-hand dominant female who reports she was working as a carpenter for Titan Construction when she fell approximately 15 feet from scaffolding on February 8, 2026. She states the scaffolding was missing guardrails and the plank shifted beneath her. She denies any loss of consciousness.

She reports her current symptoms include:
\begin{itemize}[nosep]
\item Intermittent low back pain, rated 3/10 on average
\item Occasional left leg tingling with prolonged sitting
\item Right wrist stiffness, particularly in the morning
\item Reduced grip strength in right hand
\item No current headaches or cognitive complaints
\end{itemize}

\section*{PHYSICAL EXAMINATION}

\textbf{General:} Well-developed, well-nourished female in no acute distress. Alert and cooperative.

\textbf{Right Wrist/Hand:}
\begin{itemize}[nosep]
\item Healed surgical scar from ER laceration repair
\item Mild residual swelling over distal radius
\item ROM: Flexion 55° (normal 80°), Extension 50° (normal 70°)
\item Grip strength: 18 kg (left 32 kg) - 56\% of contralateral
\item Sensation intact
\item No instability
\end{itemize}

\textbf{Lumbar Spine:}
\begin{itemize}[nosep]
\item Gait: Normal, no antalgic pattern
\item ROM: Flexion 70°, Extension 25°, Lateral flexion 30° bilaterally
\item Palpation: Mild tenderness L4-L5 paraspinals
\item SLR: Negative bilaterally
\item Neurological: Motor 5/5, Sensation intact, Reflexes 2+ symmetric
\end{itemize}

\textbf{Cognitive:} Alert, oriented, appropriate. No evidence of ongoing post-concussion syndrome.

\section*{DIAGNOSES}

\begin{enumerate}
\item Healed right distal radius fracture with residual stiffness
\item Resolved concussion
\item L4-L5 disc herniation, improving, with resolved radiculopathy
\item Lumbar strain/sprain, resolved
\end{enumerate}

\section*{DISCUSSION AND OPINIONS}

\subsection*{1. Causation}

It is my opinion, to a reasonable degree of medical certainty, that the right wrist fracture and concussion are \textbf{directly caused} by the workplace fall of February 8, 2026.

Regarding the lumbar disc herniation at L4-L5, the mechanism of injury (15-foot fall onto concrete) is consistent with acute disc injury. While the MRI noted ``mild degenerative changes,'' which may have been pre-existing, the symptomatic disc herniation with radiculopathy is \textbf{more likely than not causally related} to the workplace injury, or at minimum, the injury caused a significant aggravation of any pre-existing condition.

\subsection*{2. Maximum Medical Improvement}

The concussion has resolved. The patient is at MMI for this condition.

The right wrist fracture has healed, but the patient has residual stiffness and weakness. Given that she is still in occupational therapy, I would estimate she will reach MMI for the wrist in approximately \textbf{4-6 weeks}.

The lumbar spine has shown significant improvement with physical therapy. The radiculopathy has resolved. She may reach MMI for the lumbar condition in approximately \textbf{6-8 weeks}.

\subsection*{3. Permanent Impairment}

Using the AMA Guides to the Evaluation of Permanent Impairment, 5th Edition:

\begin{tabular}{@{}p{6cm}r@{}}
\toprule
\textbf{Condition} & \textbf{Estimated WPI} \\
\midrule
Right Wrist (if current deficits persist) & 5-7\% \\
Lumbar Spine (DRE Category II) & 5-8\% \\
Concussion & 0\% \\
\midrule
\textbf{Combined Estimate} & \textbf{10-14\%} \\
\bottomrule
\end{tabular}

\textit{Note: Final impairment rating should be deferred until MMI is reached.}

\subsection*{4. Current Treatment}

The physical therapy and occupational therapy are \textbf{reasonable and necessary}. The treatment frequency and duration are appropriate for her injuries.

\subsection*{5. Work Capacity}

Currently, Ms. Santos can perform \textbf{sedentary to light duty work} with the following restrictions:
\begin{itemize}[nosep]
\item No lifting over 15 lbs with right hand
\item No repetitive gripping with right hand
\item Ability to change positions every 30-45 minutes
\item No climbing ladders or scaffolding
\item No work at heights
\end{itemize}

I anticipate she \textbf{can return to full duty} construction work in approximately \textbf{8-12 weeks} if she continues to progress as expected.

\subsection*{6. Future Medical Care}

\begin{itemize}[nosep]
\item Complete current course of physical therapy (6-8 more weeks)
\item Complete current course of occupational therapy (4-6 more weeks)
\item Possible need for 1-2 cortisone injections if symptoms plateau
\item Low probability of requiring lumbar surgery given current improvement
\item Annual follow-up orthopedic visits for 1-2 years
\end{itemize}

\section*{SUMMARY}

Ms. Santos sustained legitimate injuries from her February 8, 2026 workplace fall including a right wrist fracture, concussion, and lumbar disc herniation. She has made good progress with treatment. She is not yet at MMI but is expected to reach MMI in approximately 2-3 months. Anticipated permanent impairment is in the range of 10-14\% whole person. She should be able to return to full duty work within 8-12 weeks.

\vspace{2em}

I declare under penalty of perjury that the information contained in this report is true and correct to the best of my knowledge.

\vspace{2em}

\begin{tabular}{@{}p{7cm}p{7cm}@{}}
\hrulefill & \\
\textbf{Richard Stevenson, MD} & \textbf{Date:} April 15, 2026 \\
Board Certified Orthopedic Surgery & \\
CA Medical License \#A-XXXXX & \\
\end{tabular}

\end{document}
